\section{Coordinate bridges and opposed paths}
\label{sec:bridges}

By Proposition~\ref{prop:coordinate-reduction}, it remains to exclude
coexistence in a coordinate direction. We therefore assume for the rest of
the proof that both $A_{e_1}$ and $A_{-e_1}$ have positive probability.
The two corresponding slab laws have finite mean width and finite mean
radius, by Propositions~\ref{prop:mean-width} and
\ref{prop:radius-moment}. We will arrange independent sequences of these
slabs in opposite chronological directions. The contradiction will come
from two estimates on how often the resulting paths approach each other.
The bridge and opposed-path constructions follow the corresponding iid
strategy in \cite[Sections~4.2--4.3]{OpenAI2026}. We give their marked
versions here, including the separation arguments for the present
environment law.

For $\sigma\in\{+,-\}$, with $\sigma e_1$ denoting $e_1$ or $-e_1$, write
\[
 \widehat P_\sigma=\widehat P_{\sigma e_1},\qquad
 D_\sigma=D_{\sigma e_1},\qquad p_\sigma=p_{\sigma e_1}>0,
 \qquad \nu_\sigma=\nu_{\sigma e_1}.
\]
For a walk starting at the origin, let $T_j^\sigma$ be the first hit of
signed coordinate height $\sigma x_1=j$. For integers $n\ge0$, set
\begin{equation}
 u_n^\sigma=P_0(T_n^\sigma<T_{-1}^\sigma),
 \qquad F_\sigma(n)=\nu_\sigma(L>n).
 \label{eq:hitting-width-tail}
\end{equation}
Thus $u_0^\sigma=1$, and $u_n^\sigma$ decreases to a limit at least
$p_\sigma$: on $D_\sigma$, the signed height reaches every positive integer
almost surely. Constants below may depend on the fixed environment law,
$d,\kappa,R$, and the chosen script, but not on any of the scale parameters.

\subsection{Bridges ending at prescribed cuts}

Under $\widehat P_\sigma$, the widths of successive slabs form a renewal
process on the nonnegative integers. Let $\bar u_n^\sigma$ be its renewal
mass, including the initial renewal at zero. Equivalently, it is the
probability of a cut at signed height $n$, with the initial boundary counted
as a cut. To identify the corresponding finite path law, for $H\ge M$ let
$\mathcal D_H^\sigma$ be the following event for the raw marked walk: it
first hits $H-M$ before $-1$ in signed height, and from that hit it takes
the prescribed $M$-step script. The script ends at its first hit of $H$.

If $\bar u_H^\sigma>0$, let $\mathcal B_H^\sigma$ denote the law of the slab
list through height $H$, conditional on a cut there. At $H=0$ this law is
concentrated on the empty list. We call $\mathcal B_H^\sigma$ the
\emph{exact-cut bridge law}.

\begin{proposition}[Exact-cut bridges]
\label{prop:exact-cut-bridges}
The renewal masses satisfy
\begin{equation}
 \bar u_0^\sigma=1,\qquad
 \bar u_j^\sigma=0\quad(1\le j<M),\qquad
 \bar u_H^\sigma=\xi u_{H-M}^\sigma\quad(H\ge M).
 \label{eq:bridge-renewal}
\end{equation}
For a list of slabs $(\gamma_1,\ldots,\gamma_m)$ with total width $H$, its
mass under $\mathcal B_H^\sigma$ is
\begin{equation}
 \mathcal B_H^\sigma(\gamma_1,\ldots,\gamma_m)
   =\frac{\prod_{i=1}^m W(\gamma_i)}{\bar u_H^\sigma}.
 \label{eq:bridge-law}
\end{equation}
In particular, reversing the list order preserves this law; the steps
within each word retain their original order.

For $H\ge M$, concatenating the list gives the raw marked path through its
first hit of $H$, conditioned on $\mathcal D_H^\sigma$. Each such path has
a unique parsing into slabs. Concatenating any prescribed slab words, or
successive bridges of prescribed widths, gives a raw path weight equal to
the product of the constituent raw weights.
\end{proposition}

\begin{proof}
The raw probability of $\mathcal D_H^\sigma$ is
$\xi u_{H-M}^\sigma$. Its terminal script begins at a record and ends at
a candidate of height $H$. By the script gap and
Lemma~\ref{lem:separated-weights}, imposing a suffix that stays at signed
height at least $H$ multiplies this probability by $p_\sigma$.
Conversely, a cut at $H$ must be the end of the script beginning at the
first hit of $H-M$: coordinate records are first hits, and each scripted
step increases the signed height by one. Dividing by $p_\sigma$ to
condition on $D_\sigma$ proves \eqref{eq:bridge-renewal}. No cut can have
width less than $M$.

The iid slab law from Proposition~\ref{prop:cut-law} now gives
\eqref{eq:bridge-law}. Its numerator and the constraint on total width are
unchanged when the list is reversed. Alternatively, for each fixed prefix
through $H$ realizing $\mathcal D_H^\sigma$, the cut suffix supplies the
same factor $p_\sigma$. Conditional on the cut at $H$ under
$\widehat P_\sigma$, the prefix therefore has mass
$W(\text{prefix})/\bar u_H^\sigma$. This is exactly its raw conditional
mass given $\mathcal D_H^\sigma$.

The parsing can be recovered from the finite prefix alone. Select the
candidate ends whose heights are not undercut later in the prefix, and
split at those ends. A suffix staying at or above $H$ cannot change this
selection, because $H$ is a strict record height. Between successive
selected candidates, every earlier candidate is undercut, which is
precisely the defining condition for a slab. A script cannot straddle a
selected boundary: it would overlap the script ending there. Thus these
are exactly the slab boundaries, and the parsing is unique.

Finally, active departures of a slab ending at signed height $b$ lie
strictly below $b-M$. All departures of the following slabs lie at or
above $b$. These row sets are more than $R$ apart in infinity distance.
Applying Lemma~\ref{lem:separated-weights} successively, and using
translation invariance, factors the weight of any prescribed
concatenation. The same argument applies after grouping the slabs into
bridges of prescribed widths. Each group boundary is reached on a first
hit, so the grouping introduces no extra path multiplicity.
\end{proof}

The distinction between $u_n^\sigma$ and $\bar u_n^\sigma$ will matter.
The former are ordinary hitting probabilities and decrease with $n$;
the latter have a gap before the first possible script end. The next
estimate retains the contribution of that gap.

\begin{lemma}[Width tails control hitting differences]
\label{lem:width-tail}
For either sign and all integers $n'\ge n\ge0$,
\begin{equation}
 u_n^\sigma-u_{n'}^\sigma
   \le C(n'-n)F_\sigma(n),
 \qquad
 \sum_{l\ge0}2^lF_\sigma(2^l)<\infty.
 \label{eq:width-tail-comparison}
\end{equation}
\end{lemma}

\begin{proof}
Fix a sign and suppress it. Partition according to the last renewal at
or before $m$. The next width exceeds the distance from that renewal to
$m$, so
\[
 \sum_{k=0}^m F(k)\bar u_{m-k}=1.
\]
Subtracting the identity at $m+1$, and using $\bar u_0=1$, gives
\[
 \sum_{k=0}^m F(k)
       (\bar u_{m-k}-\bar u_{m+1-k})=F(m+1).
\]
By \eqref{eq:bridge-renewal}, every adjacent difference in this sum is
nonnegative except the one with $m-k=M-1$, which equals $-\xi$.
Take $m=n+M$. The exceptional index is $k=n+1$, while the term $k=0$
equals $\xi(u_n-u_{n+1})$. Dropping the other nonnegative terms yields
\[
 \xi(u_n-u_{n+1})
   \le F(n+M+1)+\xi F(n+1)
   \le(1+\xi)F(n).
\]
Summation from $n$ to $n'-1$, using that $F$ decreases, proves the first
bound. The second follows by grouping the ordinary tail sum
$\sum_{k\ge0}F(k)=E_{\nu_\sigma}L<\infty$ into dyadic intervals.
\end{proof}

\subsection{Opposed tapes and common blocks}

Take two independent sequences of iid relative slab words, with laws
$\nu_+$ and $\nu_-$. We call these sequences the positive and negative
\emph{tapes}, and denote their joint law and expectation by
$\mathbf P$ and $\mathbf E$. Write
$\pi x=(x_2,\ldots,x_d)$ for the lateral coordinate of $x$.
Place the words in space as follows.
\begin{itemize}[leftmargin=*,itemsep=3pt]
\item The first positive word has its terminal site at $0$. Each
subsequent positive word has its terminal site at the starting site of
the preceding word. Thus tape order moves downward, whereas a finite
initial portion is traversed chronologically by taking its words in
reverse list order, always following the steps of each word forward.
\item The negative words are concatenated in their original order,
starting at $-e_1$. This produces a downward path with law
$\widehat P_-$ translated to $-e_1$.
\end{itemize}
Backward placement of regeneration pieces also appears in Berger's
backward-path construction \cite[Sections~2--3]{Berger2008}.
We measure \emph{depth} by $z=-x_1$. A word of width $L$, preceded by
total width $s$ on its tape, has all its departure depths in
\begin{equation}
                    (s,s+L]\cap\ZZ.
 \label{eq:departure-depths}
\end{equation}
For the positive tape this excludes the terminal arrival at depth $s$.
For the negative tape the same interval results from the initial
one-level offset: its start has depth $s+1$, and its terminal arrival has
depth $s+L+1$.

A \emph{contact} is a positive departure site at infinity distance at
most $R$ from some negative departure site. All departures count,
regardless of their marks. The depth and, subsequently, the block index
of a contact always refer to the positive departure. There is at least
one contact: the last positive departure before its terminal site $0$
is $-e_1$, the initial site of the negative path.

To compare contacts at different depths, we group the two tapes using
their common width renewals. Starting from width zero, end the first
common block at the first positive integer that is a width sum on both
tapes; continue with successive common width sums. The next proposition
shows that this grouping gives iid finite pieces with the spatial
integrability needed below. This is the intersection construction for
independent renewal processes; see \cite{AlexanderBerger2016}.

\begin{proposition}[Common blocks]
\label{prop:common-blocks}
The common blocks exist indefinitely, and their pairs of slab lists are
iid. Their widths $K_i$ have a common law satisfying
\[
                  M\le K_i<\infty\quad\text{a.s.},
             \qquad \mu:=\mathbf E K_i<\infty.
\]
The expected sum of the radii of all slabs on both sides of one common
block is finite. In particular, if
\begin{equation}
 W'_0=0,\qquad W'_i=\sum_{j=1}^iK_j,
 \label{eq:common-width-sums}
\end{equation}
then $W'_i/i\to\mu$ almost surely.
\end{proposition}

\begin{proof}
Let $K$ be the first positive common width, allowing $K=\infty$ for now.
For each fixed pair of finite prefixes that realize $K=k$, the fact
that $k$ is their first common width is decided by those prefixes. Their
unused tails are independent tapes with the initial laws. Enumerating
these prefix pairs therefore gives the renewal recursion
\[
 c_0=1,\qquad
 c_n=\sum_{k=1}^n\mathbf P(K=k)c_{n-k}\quad(n\ge1),
 \qquad c_n=\bar u_n^+\bar u_n^-.
\]
Set $f(t)=\sum_{k\ge1}\mathbf P(K=k)t^k$ for $0<t<1$. Nonnegative
power-series summation gives
\[
 \sum_{n\ge0}c_nt^n=\frac{1}{1-f(t)}.
\]
Equation~\eqref{eq:bridge-renewal} implies
$c_n\ge\xi^2p_+p_->0$ for every $n\ge M$. Hence
$\sum c_nt^n\ge\xi^2p_+p_-t^M/(1-t)$, and consequently
$(1-f(t))/(1-t)$ stays bounded as $t\uparrow1$. First this forces
$f(t)\to1$, proving $K<\infty$ almost surely. We can then write
\[
 \frac{1-f(t)}{1-t}
  =\mathbf E\frac{1-t^K}{1-t}\ \uparrow\ \mathbf E K,
\]
which proves finite mean. The same prefix factorization at each common
renewal proves that the successive pairs of lists are iid. The lower
bound $K\ge M$ follows from the minimum slab width. The strong law gives
the assertion about $W'_i$.

Let $N_+$ be the number of positive words in the first common block.
Since every width is at least one, $N_+\le K$. The event $\{N_+\ge i\}$
is determined by the first $i-1$ positive widths and the entire negative
tape: none of those positive width sums may be a negative width sum.
It is therefore independent of the $i$th positive word and its radius
$\mathcal R_i^+$. Tonelli's theorem gives
\[
 \mathbf E\sum_{i=1}^{N_+}\mathcal R_i^+
  =\sum_{i\ge1}\mathbf P(N_+\ge i)E_{\nu_+}\mathcal R
  =\mathbf E N_+\,E_{\nu_+}\mathcal R<\infty.
\]
The identical argument for the negative tape proves the claimed
integrability, using Proposition~\ref{prop:radius-moment}.
\end{proof}

The notation $W'_i$ records cumulative width and is distinct from the
word weight $W(\gamma)$. By \eqref{eq:departure-depths}, departures in
common block $i$ on either tape have depths in
$(W'_{i-1},W'_i]$. A contact in positive block $i$ can involve a negative
departure only in block $i-1$, $i$, or $i+1$: skipping an entire common
block would give a depth difference greater than $M>R$.

For each common block let $S'_i\in\ZZ^{d-1}$ be the sum of the lateral
displacements of all its positive and negative words, with both
displacements measured in the words' original chronological directions.
The $S'_i$ are iid and
\begin{equation}
                    \mathbf E\abs{S'_i}<\infty.
 \label{eq:lateral-moment}
\end{equation}
Indeed, their norms are bounded by the sums of radii controlled in
Proposition~\ref{prop:common-blocks}. As we pass downward through a
block, the negative anchor moves by its chronological displacement,
whereas the positive anchor moves by the negative of its chronological
displacement. The lateral gap, negative anchor minus positive anchor,
therefore increases by $S'_i$.

We now have both a common block index for the two tapes and integrable
iid increments for their lateral gap. The quantity on which the two
remaining estimates will disagree is
\begin{equation}
 m(J)=\sum_{j=1}^J\mathbf P\bigl\{
  \text{a contact occurs in a positive block of index in }
                      (2^j,2^{j+1}]\bigr\},\qquad J\ge1.
 \label{eq:contact-count}
\end{equation}
Thus $m(J)$ is the expected number of these dyadic block-index intervals
that contain a contact. Section~\ref{sec:entropy} first gives an upper
bound by comparing the dependent placement of the two tapes with
independent bridges.
